\documentclass{amsart}

\usepackage[utf8]{inputenc}
\usepackage{tikzit}

\input{sample.tikzstyles}

\title{ISIS1105\\ Parcial 1 - 199320}


\begin{document}
\maketitle

\section{Puntos fijos de un arreglo}
Sea $A$ un arreglo de longitud $n$ cuyas entradas son números enteros y tal que
\[
	A[0] < A[1] < \cdots < A[n-1].
\]
Queremos resolver el problema de encontrar un índice $i$ tal que $A[i]= i$, si es
que existe.

\subsection{[5 Pts]} Describa las entradas y salidas del problema.

\begin{tabular}{|| l| l| l| l||}
	\hline
	{\bf E/S} & {\bf Nombre} & {\bf Tipo} & {\bf Descripción}\\
	\hline
	        \hspace{1cm} & \hspace{3cm} & \hspace{2cm} & \hspace{4cm} \\
	          &          &          & \\
	          &          &          & \\
	          &          &          & \\
	          &          &          & \\
	          &          &          & \\
	          &          &          & \\
	          &          &          & \\
	          &          &          & \\
	          &          &          & \\
	          &          &          & \\
		  \hline
\end{tabular}

\subsection{[8 Pts]} Escriba un algoritmo recursivo que resuelva el problema en
tiempo $O(\log n)$.

\subsection{[7 Pts]} Defina una ecuación de recurrencia que exprese la complejidad
temporal de su solución y resuélvala.

\newpage
\section{Juego de carnaval}
Suponga que en la galería de tiro de un carnaval hay dos opciones para ganar puntos:
\begin{description}
	\item [Opción 1] Golpear un solo blanco, obteniendo el valor inscrito.
		\[
			\begin{array}{c}
				\tikzfig{hitone}\\
				\mbox{\emph{Puntaje $+2$}}
			\end{array}
		\]
		
	\item [Opción 2] Golpear el disco que conecta dos consecutivos, obteniendo
		el producto de los valores inscritos.
		\[
			\begin{array}{c}
				\tikzfig{hittwo}\\ 
				\mbox{\emph{Puntaje $+18$}}
			\end{array}
		\]
	\item [Tenga en cuenta] El puntaje inicial es cero y pueden haber blancos
		con puntajes negativos. Está permitido no golpear un blanco para
		evitar perder puntos.
\end{description}

\newpage
\subsection{[10 Pts]} Describa las entradas y salidas del problema.

\begin{tabular}{|| l| l| l| l||}
	\hline
	{\bf E/S} & {\bf Nombre} & {\bf Tipo} & {\bf Descripción}\\
	\hline
	        \hspace{1cm} & \hspace{3cm} & \hspace{2cm} & \hspace{4cm} \\
	          &          &          & \\
	          &          &          & \\
	          &          &          & \\
	          &          &          & \\
	          &          &          & \\
	          &          &          & \\
	          &          &          & \\
	          &          &          & \\
	          &          &          & \\
	          &          &          & \\
		  \hline
\end{tabular}

\subsection{[20 Pts]} Defina una ecuación de recurrencia para $L(k)$, el máximo
puntaje que se puede obtener jugando con los blancos $0, \ldots, k-1$.

\subsection{[10 Pts]} Dibujar el grafo de necesidades asociado con esta ecuación.
Explique qué representan los vértices y qué representan las aristas del grafo.

\subsection{[20 Pts]} Desarrolle un algoritmo de programación dinámica para
resolver este problema.

\newpage
\section{Divisibilidad}
Considere el grafo dirigido $G_N$ cuyos nodos son los enteros \emph{positivos}
menores o iguales a $N$ y donde la arista $a \to b$ pertenece al grafo si y sólo
si $b= ap$ donde $p$ es un número primo. Abajo encuentra un dibujo de $G_{10}$.

\ctikzfig{div10}

Dado $G_N$ y dos enteros positivos distintos $x, y$ menores o iguales a $N$,
queremos determinar si alguno es múltiplo del otro.

\subsection{[5 Pts]} Describa las entradas y salidas del problema.

\begin{tabular}{|| l| l| l| l||}
	\hline
	{\bf E/S} & {\bf Nombre} & {\bf Tipo} & {\bf Descripción}\\
	\hline
	        \hspace{1cm} & \hspace{3cm} & \hspace{2cm} & \hspace{4cm} \\
	          &          &          & \\
	          &          &          & \\
	          &          &          & \\
	          &          &          & \\
	          &          &          & \\
	          &          &          & \\
	          &          &          & \\
	          &          &          & \\
	          &          &          & \\
	          &          &          & \\
		  \hline
\end{tabular}

\subsection{[10 Pts]} Determine de qué problema general de grafos este ejercicio
es una instancia particular.

\subsection{[10 Pts]} Escriba un algoritmo para resolver el problema.

\end{document}
